\begin{quote}
\emph{Lane A constructs the sharp local operator algebra and verifies the OS axioms for the limiting local state.
The reflection-positivity input is taken from the unflowed Wilson state (Section~\ref{sec:patching});
no reflection-positivity property is assumed for the auxiliary flowed observable sector.}
\end{quote}

\paragraph{OS-positive base state.}
Fix a weak-window target value $u_{\mathrm{ren}}\in(0,u_\ast]$ and let
\[
a_n:=\ell_0 2^{-n},\qquad \beta_n:=\beta(a_n;u_{\mathrm{ren}})
\]
be the associated tuned dyadic trajectory.
Theorem~\ref{thm:continuum_tightness} furnishes the bounded positive-time continuum Wilson state $\omega_{\mathrm{bd}}^{(u_{\mathrm{ren}})}$ on $\mathcal A_+$ together with the actual dyadic OS matrix elements needed downstream.
This is the only reflection-positive base state used by Lane~A.

\paragraph{Flowed probes.}
Lane~C also provides continuum Schwinger functions for the flowed probe family $\mathcal A_{\mathrm{flow}}$ (Theorem~\ref{thm:C3U}).
Because gradient flow smears in the Euclidean time direction, reflection positivity does not transfer to the full flowed algebra by a support argument,
and we do not use (or require) reflection positivity at positive flow time.
The role of flow in Lane~A is renormalization and identification: it produces UV-regular probes and a small-flow-time expansion that
isolates the sharp local operators.

\paragraph{Already-constructed Lane~C flowed state.}
For the same tuned dyadic trajectory, Theorem~\ref{thm:C3U} already furnishes a unique continuum Euclidean state
\[
\omega_{\mathrm{flow}}^{(u_{\mathrm{ren}})}:\mathcal A_{\mathrm{flow}}\to\mathbb C.
\]
Lane~A does not construct a second state on flowed observables: it keeps $\omega_{\mathrm{flow}}^{(u_{\mathrm{ren}})}$ fixed on $\mathcal A_{\mathrm{flow}}$ and extends it to sharp-local insertions by inverse-flow limits plus quantitative inverse control.
For brevity below, write $\omega_{\mathrm{flow}}:=\omega_{\mathrm{flow}}^{(u_{\mathrm{ren}})}$.

\begin{lemma}[RP-safe positive-time representatives for flowed probes]\label{lem:RP_safe_flow_rep}
Fix $\varepsilon>0$ and let $f$ be supported in $\{x_0\ge \varepsilon\}$.
For $t\in(0,t_\ast]$ and $R\ge 1$ with $R\sqrt t\le \varepsilon/2$, define a truncation operator $T_{R,t}$ on lattice configurations
by replacing the initial field outside the Euclidean neighborhood $B(\mathrm{supp}(f),R\sqrt t)$ by a fixed reference configuration.
Set
\[
\mathcal O^{(R)}_{i,t}(f):=\mathcal O_{i,t}(f)\circ T_{R,t}.
\]
Then $\mathcal O^{(R)}_{i,t}(f)$ is a bounded cylinder functional depending only on underlying variables supported in $\{x_0>0\}$ and hence lies in the unflowed positive-time algebra $\mathcal A_+$.

Moreover, for any bounded $G\in\mathcal A_+$,
\[
\big|\omega_a(\Theta G\,\mathcal O_{i,t}(f)) - \omega_a(\Theta G\,\mathcal O^{(R)}_{i,t}(f))\big|
\le \|G\|_\infty\,\|\mathcal O_{i,t}(f)-\mathcal O^{(R)}_{i,t}(f)\|_\infty.
\]
By the flow quasi-locality bounds (Theorem~\ref{thm:tail_R}), the truncation error satisfies
\[
\|\mathcal O_{i,t}(f)-\mathcal O^{(R)}_{i,t}(f)\|_\infty\ \le\ C\,t^{-5}e^{-cR^2}
\]
(on the good-flow chart regime; Lemma~\ref{lem:bad_event_irrelevant} upgrades this to an unconditional bound in OS/RP pairings).
Choosing $R=R(t)\asymp\sqrt{\log(1/t)}$ as in Corollary~\ref{cor:R_choice} yields an error $O(t^\eta)$ for any prescribed $\eta>0$.

Consequently, every use of reflection positivity and OS inner products in this submission is applied only to elements of $\mathcal A_+$:
whenever a flowed probe appears inside an OS/RP expression, it is understood to be replaced by a positive-time representative
$\mathcal O^{(R(t))}_{i,t}(f)\in\mathcal A_+$ with a negligible error at the relevant scale.
\end{lemma}

\begin{proof}
The inclusion $B(\mathrm{supp}(f),R\sqrt t)\subset\{x_0>0\}$ follows from $R\sqrt t\le \varepsilon/2$.
The expectation difference bound is the trivial $\|G\|_\infty$ estimate.
The stated truncation error is exactly the $L^\infty$ consequence of Theorem~\ref{thm:tail_R}; the power-saving choice of $R(t)$ is Corollary~\ref{cor:R_choice}.
\end{proof}

\begin{lemma}[Bad-event contribution is negligible in OS/RP pairings]\label{lem:bad_event_irrelevant}
Fix $\varepsilon>0$ and let $f$ be supported in $\{x_0\ge\varepsilon\}$.
Let $t\in(0,t_\ast]$ and $R\ge 1$ satisfy $R\sqrt t\le\varepsilon/2$, and let $\mathcal O^{(R)}_{i,t}(f)$ be the positive-time representative from
Lemma~\ref{lem:RP_safe_flow_rep}.
Let $\Lambda$ be the finite plaquette set in the neighborhood relevant for $f$ and set
\[
G_{\Lambda}(\varepsilon):=\bigcap_{p\in\Lambda}\{\theta(U_p)\le\varepsilon\}
\]
(as in Lemma~\ref{lem:local_tail}).
Then for any bounded $G\in\mathcal A_+$,
\begin{align}\label{eq:bad_event_irrelevant_bound}
\big|\omega_a(\Theta G\,\mathcal O_{i,t}(f)) - \omega_a(\Theta G\,\mathcal O^{(R)}_{i,t}(f))\big|
&\le \|G\|_\infty\Big( C\,t^{-5}e^{-cR^2} + 2\,\|\mathcal O_{i,t}(f)\|_\infty\,\mathbb P_{\omega_a}\big(G_{\Lambda}(\varepsilon)^c\big)\Big).
\end{align}
Moreover, Lemma~\ref{lem:local_tail} gives the explicit bound
\[
\mathbb P_{\omega_a}\big(G_{\Lambda}(\varepsilon)^c\big)
\le |\Lambda|\,C_{\mathrm{tail}}(G)\,\varepsilon^{-2d_G/3}\,\exp\big(-c_{\mathrm{tail}}(G)\,\beta\,\varepsilon^2\big),
\qquad d_G:=\dim(G),
\]
valid for $\varepsilon\le \varepsilon_{\mathrm{tail}}(G)$ as in Lemma~\ref{lem:local_tail}.
In the AF window, $u\asymp \beta^{-1}$ along the tuned trajectory, so the stretched-exponential tail is \emph{superpolynomial} in $u$:
for every $M$ there exists $u_0(M)>0$ such that $u\le u_0(M)$ implies
\begin{equation}\label{eq:bad_event_prob_superpoly}
\mathbb P_{\omega_a}\big(G_{\Lambda}(\varepsilon)^c\big)\le u^{2M}.
\end{equation}
Consequently, choosing $R=R(t)\asymp\sqrt{\log(1/t)}$ yields a power-saving $O(t^{\eta})$ truncation error from the good-event term.
For the bad-event term, fix $t\in(0,t_\ast]$ and $\eta>0$, choose $M$ large, and then (along the tuned trajectory) take the continuum limit $a\downarrow 0$
far enough into the AF window so that the corresponding weak coupling satisfies $u(a)^M\le t^{\eta}$.
This makes the second term in \eqref{eq:bad_event_irrelevant_bound} also $O(t^{\eta})$ at that fixed $t$.
In particular, all RP/OS pairings may be evaluated using the positive-time representatives with an unconditional $O(t^{\eta})$ error,
and the subsequent SFTE argument then sends $t\downarrow 0$.
\end{lemma}

\begin{proof}
Decompose the expectation difference into contributions from $G_{\Lambda}(\varepsilon)$ and its complement.
On $G_{\Lambda}(\varepsilon)$, Theorem~\ref{thm:tail_R} gives the $L^\infty$ truncation bound
$\|\mathcal O_{i,t}(f)-\mathcal O^{(R)}_{i,t}(f)\|_\infty\le C t^{-5}e^{-cR^2}$.
On $G_{\Lambda}(\varepsilon)^c$, use the trivial bound
$|\omega_a(\Theta G\,\mathcal O_{i,t}(f)) - \omega_a(\Theta G\,\mathcal O^{(R)}_{i,t}(f))|
\le 2\|G\|_\infty\,\|\mathcal O_{i,t}(f)\|_\infty\,\mathbb P_{\omega_a}(G_{\Lambda}(\varepsilon)^c)$.
This yields \eqref{eq:bad_event_irrelevant_bound}.
The superpolynomial probability estimate \eqref{eq:bad_event_prob_superpoly} follows from Lemma~\ref{lem:local_tail} and the elementary calculus fact that $e^{-c/u}\le u^{2M}$ for all sufficiently small $u$ (with $c>0$ fixed).
\end{proof}

\paragraph{Small-flow-time expansion and inverse mixing.}
Fix a finite sector-basis of flowed probes $\{\mathcal O_{i,t}(f)\}_{i=1}^{N}$ and a corresponding sector-basis of renormalized local operators
$\{\mathcal O^{\mathrm{ren}}_{j}(f;\mu)\}_{j=1}^{N}$.
The small-flow-time expansion (SFTE) asserts, for $t\downarrow 0$,
\begin{equation}\label{eq:SFTE_matrix_LaneA}
\mathcal O_{i,t}(f)
=\sum_{j=1}^{N} C_{ij}(t;\mu)\,\mathcal O^{\mathrm{ren}}_{j}(f;\mu)
+\mathcal R_{i}(t;f;\mu),
\end{equation}
in the sense of mixed Schwinger functions against arbitrary finite collections of renormalized local insertions.
For the closure theorem below we do \emph{not} use bare invertibility of $C(t;\mu)$ alone.
Instead the abstract inverse-flow mechanism is that for the sector family under consideration there exist $t_0\in(0,t_\ast]$ and a scale-dependent quantitative inverse control function
$\Lambda_{\mathrm{inv}}(t;\mu)\ge 1$ such that
\begin{equation}\label{eq:LaneA_quant_inverse_bound}
\max_{1\le i,j\le N}\big|\big(C(t;\mu)^{-1}\big)_{ij}\big|
\le M_{\mathrm{inv}}(\mu)\,\Lambda_{\mathrm{inv}}(t;\mu)
\qquad (0<t\le t_0),
\end{equation}
and such that the remainders become negligible after multiplication by the same inverse factor in all mixed correlator channels relevant to the target local state:
\begin{equation}\label{eq:LaneA_inverse_remainder_compat}
\Lambda_{\mathrm{inv}}(t;\mu)\,
\Big|
\omega\Big(
\big(\mathcal R_i(t;f;\mu)-\omega(\mathcal R_i(t;f;\mu))\big)
\prod_{\ell=1}^k \mathcal X_\ell
\Big)
\Big|
\xrightarrow[t\downarrow 0]{}0.
\end{equation}
Here $\omega$ is only a placeholder for the eventual sharp-local extension functional in this abstract mechanism.
The special case $\Lambda_{\mathrm{inv}}\equiv 1$ recovers the uniformly bounded inverse situation, but the theorem is designed to accommodate the
polylogarithmic inverse growth established later in Section~\ref{sec:sharp_local}.
For the public Lane~A theorem below, however, no provisional extension functional appears in the hypotheses:
Theorem~\ref{thm:finite_truncation_inverse_control} and Corollary~\ref{cor:finite_truncation_SFTE} supply the capwise inverse bound and the remainder compatibility intrinsically on every finite engineering cap.
In particular, Corollary~\ref{cor:finite_truncation_SFTE} rewrites the abstract condition \eqref{eq:LaneA_inverse_remainder_compat}
as a concrete family of capwise OS/core matrix-element bounds, so the proof of Theorem~\ref{thm:LaneA_closure}
never invokes an extension functional before the finite-cap state has been defined.

Define the (centered) inverse-flow approximants
\begin{equation}\label{eq:inverse_flow_approx}
\mathcal O^{\mathrm{app}}_{j}(t;f;\mu)
:=\sum_{i=1}^{N}\big(C(t;\mu)^{-1}\big)_{ji}\,\big(\mathcal O_{i,t}(f)-\omega_{\mathrm{flow}}(\mathcal O_{i,t}(f))\big),
\end{equation}
so the centering in the inverse-flow approximants is taken with respect to the already-constructed Lane~C flowed state.

\paragraph{Lane A extension theorem.}
The next statement is the intrinsic finite-cap Lane~A output needed downstream.

\paragraph{How to read the full-scope statement.}
There are three closure levels, and separating them makes the widened Lane~A theorem much easier to audit.
\begin{enumerate}[label=(\roman*),leftmargin=2.3em]
\item \emph{Finite audited sectors.} The scalar/tensor/loop examples show concretely how inverse control is verified, but after the quotient-block hardening they are illustrations rather than the source of the general breadth claim.
\item \emph{Finite engineering truncations.} Corollary~\ref{cor:LaneA_all_finite_truncations} is the breadth input needed here:
it says that both the quantitative inverse-control package and the SFTE/remainder-compatibility package hold intrinsically on every finite cap of the canonical sharp-local basis.
Its blockwise input is now theorem-level: Appendix~\ref{app:normalization_crosscheck} proves the coefficient-extraction theorem for arbitrary exact-dimension quotient blocks,
so the canonical diagonal UV problem is a literal finite-dimensional statement before the truncation assembly is invoked.
\item \emph{Inductive union.} Corollary~\ref{cor:LaneA_full_scope} then uses the coherence of the normalization scheme to identify the same renormalized fields inside every larger cap and passes to the full sharp-local algebra.
\end{enumerate}
Thus the theorem below should be read as the finite-cap Lane~A extension mechanism with no extra sectorwise assumptions,
while Corollary~\ref{cor:LaneA_full_scope} is the separate inductive-union upgrade to the full sharp-local algebra.

\begin{theorem}[Lane~A finite-cap extension of the Lane~C flowed state]\label{thm:LaneA_closure}
Fix a weak-window target value $u_{\mathrm{ren}}\in(0,u_\ast]$ and let
\[
a_n:=\ell_0 2^{-n},\qquad \beta_n:=\beta(a_n;u_{\mathrm{ren}})
\]
be the tuned dyadic trajectory.
Let $\omega_{\mathrm{flow}}^{(u_{\mathrm{ren}})}$ be the already-constructed continuum Euclidean state on $\mathcal A_{\mathrm{flow}}$ furnished by Theorem~\ref{thm:C3U} along this same tuning, and let
$\omega_{\mathrm{bd}}^{(u_{\mathrm{ren}})}$ be the bounded positive-time continuum Wilson state on $\mathcal A_+$.
Let
\[
\mathcal K_{u_{\mathrm{ren}}}(F,G;s):=\lim_{n\to\infty}\omega_{a_n}(\Theta F\,\tau_s G)
\qquad(F,G\in\mathcal A_+,\ s\in\mathbb D_{\mathrm{dyad}})
\]
be the dyadic OS matrix-element limits furnished by Theorem~\ref{thm:continuum_tightness}.

Fix an engineering cutoff $\Delta_{\max}\ge 0$ and let
\[
\mathcal A_{\mathrm{flow}}^{(\le \Delta_{\max})}\subset\mathcal A_{\mathrm{flow}}
\]
denote the capwise bounded flowed subalgebra generated by the canonical flowed lifts entering that truncation.
Let $\mathfrak A_{\mathrm{loc}}^{(\le \Delta_{\max})}$ be the $\ast$-algebra generated by $\mathcal A_{\mathrm{flow}}^{(\le \Delta_{\max})}$, $\mathcal A_+$, and all canonical renormalized local generators of engineering dimension $\le \Delta_{\max}$.
Let
\[
\mathbf Y_t^{(\le \Delta_{\max})}(f)
=
C^{(\le \Delta_{\max})}(t;\mu)\,\mathbf O_{\mathrm{ren}}^{(\le \Delta_{\max})}(f;\mu)
+\mathbf R_t^{(\le \Delta_{\max})}(f;\mu)
\]
be the canonical finite-cap small-flow-time expansion furnished by Corollary~\ref{cor:finite_truncation_SFTE}, written in the capwise bases
\[
\begin{aligned}
\mathbf Y_t^{(\le \Delta_{\max})}(f)
&=\bigl(Y_{1,t}^{(\le \Delta_{\max})}(f),\dots,Y_{N,t}^{(\le \Delta_{\max})}(f)\bigr)^{\mathsf T},\\
\mathbf O_{\mathrm{ren}}^{(\le \Delta_{\max})}(f;\mu)
&=\bigl(O^{(1),\mathrm{ren}}_{\le \Delta_{\max}}(f;\mu),\dots,O^{(N),\mathrm{ren}}_{\le \Delta_{\max}}(f;\mu)\bigr)^{\mathsf T}.
\end{aligned}
\]
For each capwise basis index $j$, define the centered inverse-flow approximant
\[
\mathcal O^{\mathrm{app},(\le \Delta_{\max})}_{j}(t;f;\mu)
:=
\sum_i \big((C^{(\le \Delta_{\max})}(t;\mu))^{-1}\big)_{ji}\,
\big(Y^{(\le \Delta_{\max})}_{i,t}(f)-\omega_{\mathrm{flow}}^{(u_{\mathrm{ren}})}(Y^{(\le \Delta_{\max})}_{i,t}(f))\big).
\]
Then Theorem~\ref{thm:finite_cap_flowed_sector_bridge}, together with
Theorem~\ref{thm:finite_truncation_inverse_control} and Corollary~\ref{cor:finite_truncation_SFTE}, defines a unique sharp-local continuum state
\[
\omega_{\le \Delta_{\max}}^{(u_{\mathrm{ren}})}
\]
on $\mathfrak A_{\mathrm{loc}}^{(\le \Delta_{\max})}$ extending the restriction of $\omega_{\mathrm{flow}}^{(u_{\mathrm{ren}})}$ to $\mathcal A_{\mathrm{flow}}^{(\le \Delta_{\max})}$.
Equivalently, once the state has been constructed, the formal centered generators of Theorem~\ref{thm:finite_cap_flowed_sector_bridge} are written as
\[
\widehat O_j^{(\le \Delta_{\max})}(f;\mu)
=
O^{(j),\mathrm{ren}}_{\le \Delta_{\max}}(f;\mu)
-\omega_{\le \Delta_{\max}}^{(u_{\mathrm{ren}})}\big(O^{(j),\mathrm{ren}}_{\le \Delta_{\max}}(f;\mu)\big).
\]
Then:
\begin{enumerate}[label=(\alph*)]
\item (\textbf{Sharp local identification}) For each capwise basis index $j$ and test function $f$, the inverse-flow approximants satisfy
\[
\begin{aligned}
\lim_{t\downarrow 0}\, \omega_{\le \Delta_{\max}}^{(u_{\mathrm{ren}})}\Big(
&\mathcal O^{\mathrm{app},(\le \Delta_{\max})}_{j}(t;f;\mu)
\prod_{\ell=1}^{k} \mathcal X_{\ell}\Big)\\
={}&\omega_{\le \Delta_{\max}}^{(u_{\mathrm{ren}})}\Big(
\big(O^{(j),\mathrm{ren}}_{\le \Delta_{\max}}(f;\mu)
-\omega_{\le \Delta_{\max}}^{(u_{\mathrm{ren}})}(O^{(j),\mathrm{ren}}_{\le \Delta_{\max}}(f;\mu))\big)
\prod_{\ell=1}^{k} \mathcal X_{\ell}\Big).
\end{aligned}
\]
for every finite collection of renormalized local insertions $\{\mathcal X_{\ell}\}$ drawn from the same cap.
\item (\textbf{OS axioms for the finite-cap local state}) The Schwinger functions of these renormalized local fields satisfy the Osterwalder--Schrader axioms.
In particular, OS reflection positivity holds on the capwise local positive-time algebra, and the associated field polynomials are realized on the explicit common invariant core of bounded flowed positive-time classes recorded in Lemma~\ref{lem:finite_cap_common_core}.
\item (\textbf{Compatibility with the Lane~C flowed state and the tuned-dyadic bounded base state})
The restriction of $\omega_{\le \Delta_{\max}}^{(u_{\mathrm{ren}})}$ to the capwise flowed subalgebra equals the already-constructed Lane~C state:
\[
\omega_{\le \Delta_{\max}}^{(u_{\mathrm{ren}})}\!\restriction_{\mathcal A_{\mathrm{flow}}^{(\le \Delta_{\max})}}
=
\omega_{\mathrm{flow}}^{(u_{\mathrm{ren}})}\!\restriction_{\mathcal A_{\mathrm{flow}}^{(\le \Delta_{\max})}}.
\]
On the bounded positive-time algebra, the extension agrees with $\omega_{\mathrm{bd}}^{(u_{\mathrm{ren}})}$:
\[
\omega_{\le \Delta_{\max}}^{(u_{\mathrm{ren}})}\!\restriction_{\mathcal A_+}
=
\omega_{\mathrm{bd}}^{(u_{\mathrm{ren}})}.
\]
Moreover, for every bounded $F,G\in\mathcal A_+$ and every dyadic shift $s\in\mathbb D_{\mathrm{dyad}}$,
\[
\omega_{\le \Delta_{\max}}^{(u_{\mathrm{ren}})}(\Theta F\,\tau_s G)=\mathcal K_{u_{\mathrm{ren}}}(F,G;s).
\]
\end{enumerate}
\end{theorem}

\paragraph{Mixed-channel downstream closure node.}
For downstream Lane~A closure statements, the positive-time mixed-channel Cauchy clause enters only through Theorem~\ref{thm:finite_cap_flowed_sector_bridge}(i), with Corollary~\ref{cor:finite_cap_mixed_family_packaging} supplying the finite-family packaging/covariance node used there.
Corollary~\ref{cor:finite_truncation_SFTE} and Lemma~\ref{lem:finite_cap_mixed_channel_reduction} contain the precursor finite-cap reductions; downstream closure statements cite only Theorem~\ref{thm:finite_cap_flowed_sector_bridge}(i).
In the notation fixed there, $K_W$ denotes the associated support neighborhood, $C_W^{\mathrm{ins}}$ packages the fixed-word insertion dependence, and the finite truncation data record the remaining theorem-facing constants; expanding $C_W^{\mathrm{ins}}$ recovers $\|W\|_{\mathrm{bnd}}:=\prod_r \|B_r\|_\infty$ together with the normalized-word factor from Corollary~\ref{cor:ins_to_OS_bounded_word}.

\begin{proof}
Fix $\Delta_{\max}$ and abbreviate
\[
\mathfrak A_{\mathrm{loc}}^{(\le \Delta_{\max})}=:\mathcal A_{\mathrm{loc}}^{\mathrm{cap}},
\qquad
\omega_{\le \Delta_{\max}}^{(u_{\mathrm{ren}})}=:\omega^{(u_{\mathrm{ren}})}.
\]
Write
\[
Y_{i,t}(f):=Y^{(\le \Delta_{\max})}_{i,t}(f),
\qquad
O_i^{\mathrm{ren}}(f;\mu):=O^{(i),\mathrm{ren}}_{\le \Delta_{\max}}(f;\mu),
\qquad
C(t;\mu):=C^{(\le \Delta_{\max})}(t;\mu).
\]
Let
\[
\mathcal B_{\mathrm{cap}}^{(\le \Delta_{\max})}
:=
\ast\!\bigl(\mathcal A_{\mathrm{flow}}^{(\le \Delta_{\max})},\mathcal A_+\bigr)
\]
be the bounded capwise algebra used in Theorem~\ref{thm:finite_cap_flowed_sector_bridge}.

\paragraph{Intrinsic capwise definition of the extension.}
For the construction it is convenient to work first with the abstract centered capwise algebra generated by
$\mathcal B_{\mathrm{cap}}^{(\le \Delta_{\max})}$ and the formal symbols
$\widehat O_j^{(\le \Delta_{\max})}(f;\mu)$.
After fixing the finite family of mixed bounded words used in the construction, Theorem~\ref{thm:finite_cap_flowed_sector_bridge}, together with Corollary~\ref{cor:finite_cap_mixed_family_packaging}, furnishes a unique positive unital functional on that abstract algebra, characterized by the inverse-flow limits against the corresponding finite-family packaged bounded capwise states.
No universal mixed bounded capwise state is presupposed here; the bounded input is exactly that locally fixed finite-family package.
We denote that functional by $\omega^{(u_{\mathrm{ren}})}$.
Once the state has been fixed, the formal centered symbols are written as
\[
\widehat O_j^{(\le \Delta_{\max})}(f;\mu)
=
O_j^{\mathrm{ren}}(f;\mu)-\omega^{(u_{\mathrm{ren}})}\big(O_j^{\mathrm{ren}}(f;\mu)\big),
\]
which is exactly the notation used in the theorem statement.
Thus existence of the capwise state has been reduced to the theorem-level mixed-correlator construction, and uniqueness will follow from the same characterization on the generating set.

\paragraph{(a) Sharp local identification.}
Fix $j$ and $f$ and let $\{\mathcal X_\ell\}_{\ell=1}^k$ be any finite family of renormalized local insertions from the cap.
Write the product $\prod_{\ell=1}^k \mathcal X_\ell$ as the corresponding abstract centered capwise word.
Theorem~\ref{thm:finite_cap_flowed_sector_bridge}(i)--(ii) defines the capwise mixed-correlator functional precisely so that replacing the slot
$\widehat O_j^{(\le \Delta_{\max})}(f;\mu)$ by the bounded inverse-flow representative
$\mathcal O_j^{\mathrm{app},(\le \Delta_{\max})}(t;f;\mu)$ changes the expectation by $o(1)$ as $t\downarrow 0$, against every fixed capwise word in the remaining slots.
For this mixed-channel step we therefore cite only Theorem~\ref{thm:finite_cap_flowed_sector_bridge}(i), whose proof packages the upstream chain Theorem~\ref{thm:ins_to_OS}, Corollary~\ref{cor:ins_to_OS_bounded_word}, Proposition~\ref{prop:finite_cap_bounded_factor_insertion}, Theorem~\ref{thm:finite_truncation_inverse_control}, and Lemma~\ref{lem:finite_cap_mixed_channel_reduction}, together with Corollary~\ref{cor:finite_cap_mixed_family_packaging} for the finite-family packaging/covariance step; the fixed-word dependence is recorded there through $K_W$, $C_W^{\mathrm{ins}}$, and the finite truncation data.
Because the same theorem gives independence of the relative rates in all local slots, one may use a common parameter $t$ for the whole finite insertion family.
After identifying the centered symbol with
$O_j^{\mathrm{ren}}(f;\mu)-\omega^{(u_{\mathrm{ren}})}(O_j^{\mathrm{ren}}(f;\mu))$, this is exactly the convergence claimed in part~(a).
The same characterization fixes all mixed moments of the generating set and therefore gives uniqueness of the capwise extension.

\paragraph{(b) OS axioms for the finite-cap local state.}
We explain why the Schwinger functions of the renormalized local fields satisfy the OS axioms.

\emph{Euclidean invariance and symmetry.}
The finite-family packaged bounded capwise states entering Theorem~\ref{thm:finite_cap_flowed_sector_bridge} inherit the same Euclidean invariance on $\mathcal A_{\mathrm{flow}}^{(\le \Delta_{\max})}$ from Theorem~\ref{thm:C3U}, while Corollary~\ref{cor:finite_cap_mixed_family_packaging} records the translation covariance used in the mixed-channel step and the tuned Wilson limit supplies the bounded $\mathcal A_+$ identities.
The local Schwinger functions are defined from those finite-family packaged states by pointwise inverse-flow limits, so the same algebraic identities persist for the limiting local correlators.
Permutation symmetry is exact at finite lattice spacing for the bounded representatives and therefore survives the limit as well.

\emph{Reflection positivity.}
Fix $\varepsilon>0$ and consider the capwise local positive-time algebra generated by insertions $O_j^{\mathrm{ren}}(f;\mu)$ with $\mathrm{supp}(f)\subset\{x_0\ge\varepsilon\}$.
Let $F$ be a finite linear combination of monomials in such generators.
Choose $t\in(0,t_\ast]$ and $R\ge 1$ with $R\sqrt t\le \varepsilon/2$.
Define $F_{t,R}$ by replacing each local generator $O_j^{\mathrm{ren}}(f;\mu)$ appearing in $F$ by its inverse-flow approximant and then replacing each flowed probe inside that approximant by the RP-safe positive-time representative from Lemma~\ref{lem:RP_safe_flow_rep}.
By construction, $F_{t,R}\in\mathcal A_+$.
Since $\omega_{\mathrm{bd}}^{(u_{\mathrm{ren}})}$ is reflection positive on $\mathcal A_+$ by Theorem~\ref{thm:continuum_tightness}, one has
\[
\omega_{\mathrm{bd}}^{(u_{\mathrm{ren}})}(\Theta F_{t,R}\cdot F_{t,R})\ge 0.
\]
Now let $t\downarrow 0$ and choose $R=R(t)\asymp\sqrt{\log(1/t)}$ as in Lemma~\ref{lem:RP_safe_flow_rep}.
By Lemma~\ref{lem:RP_safe_flow_rep}, replacing flowed probes by their positive-time representatives changes mixed Schwinger functions by $o(1)$.
By part~(a), the inverse-flow approximants converge in every fixed capwise channel to the corresponding centered local insertions.
Hence
\[
\omega_{\mathrm{bd}}^{(u_{\mathrm{ren}})}(\Theta F_{t,R(t)}\cdot F_{t,R(t)})
\xrightarrow[t\downarrow 0]{}
\omega^{(u_{\mathrm{ren}})}(\Theta F\cdot F),
\]
and the limit is nonnegative.
Thus OS reflection positivity holds on the capwise local positive-time algebra.

\emph{Regularity and operator domain.}
Theorem~\ref{thm:field_algebra_truncation} and Lemma~\ref{lem:finite_cap_common_core} provide the explicit dense core
$\mathcal D^{(\le \Delta_{\max})}_{\mathrm{flow}}$ of bounded flowed positive-time classes on which all finite field polynomials are defined by graph-norm limits of bounded dyadic approximants.
Corollary~\ref{cor:finite_truncation_SFTE} identifies those graph-norm limits with the Lane~A inverse-flow representatives, and
Theorem~\ref{thm:finite_cap_flowed_sector_bridge}(iv) identifies the resulting positive-time OS/core matrix elements with the capwise mixed-correlator state just constructed.
Thus the finite-cap local fields are realized on the explicit common invariant core appearing in the theorem statement.
For regularity, the flowed correlators are bounded by the Lane~C moment/cumulant package together with the kernel bounds in Appendices~B--D; the capwise SFTE remainder bounds then pass those estimates to the renormalized local correlators, with only polylogarithmic inverse factors, which are harmless on a fixed cap.
Hence the present finite-cap Schwinger functions satisfy OS$_0$.

\paragraph{(c) Compatibility with the Lane~C flowed state and the tuned-dyadic bounded base state.}
The restriction to $\mathcal A_{\mathrm{flow}}^{(\le \Delta_{\max})}$ is
$\omega_{\mathrm{flow}}^{(u_{\mathrm{ren}})}\!\restriction_{\mathcal A_{\mathrm{flow}}^{(\le \Delta_{\max})}}$
by Theorem~\ref{thm:finite_cap_flowed_sector_bridge}(ii).
The same theorem also gives
\[
\omega^{(u_{\mathrm{ren}})}\!\restriction_{\mathcal A_+}=\omega_{\mathrm{bd}}^{(u_{\mathrm{ren}})}.
\]
Finally, for bounded $F,G\in\mathcal A_+$ and dyadic $s\in\mathbb D_{\mathrm{dyad}}$, Theorem~\ref{thm:continuum_tightness} already records the limiting matrix element
$\mathcal K_{u_{\mathrm{ren}}}(F,G;s)$ along the tuned dyadic trajectory, so the same identity holds for the capwise extension:
\[
\omega^{(u_{\mathrm{ren}})}(\Theta F\,\tau_s G)=\mathcal K_{u_{\mathrm{ren}}}(F,G;s).
\]
This completes the proof.
\end{proof}

\begin{corollary}[Lane~A preserves the tuned-dyadic bounded base state]\label{cor:LaneA_bounded_state_restriction}
Fix $u_{\mathrm{ren}}\in(0,u_\ast]$ and let $\omega^{(u_{\mathrm{ren}})}$ be the full sharp-local state furnished by Corollary~\ref{cor:LaneA_full_scope}.
Then
\[
\omega^{(u_{\mathrm{ren}})}\!\restriction_{\mathcal A_+}=\omega_{\mathrm{bd}}^{(u_{\mathrm{ren}})}
\]
and, for bounded $F,G\in\mathcal A_+$ and dyadic $s\in\mathbb D_{\mathrm{dyad}}$,
\[
\omega^{(u_{\mathrm{ren}})}(\Theta F\,\tau_s G)=\mathcal K_{u_{\mathrm{ren}}}(F,G;s).
\]
\end{corollary}

\begin{proof}
Part~(c) of Theorem~\ref{thm:LaneA_closure} gives the stated restriction and dyadic matrix-element identity on every finite engineering cap.
Corollary~\ref{cor:LaneA_full_scope} passes those compatible identities to the inductive union.
\end{proof}

\begin{remark}[Phase-II seam note for Corollaries~\ref{cor:LaneA_bounded_state_restriction}--\ref{cor:LaneA_full_scope}]
Corollary~\ref{cor:LaneA_bounded_state_restriction} is the bounded-state compatibility theorem only.
The finite-cap sharp-local OS structure belongs to Theorem~\ref{thm:LaneA_closure}(b), while the actual passage to the full sharp-local inductive union belongs only to Corollary~\ref{cor:LaneA_full_scope}.
\end{remark}


\begin{corollary}[Full sharp-local scope by finite truncations and inductive union]\label{cor:LaneA_full_scope}
For each $\Delta_{\max}\ge 0$, let $\mathfrak A_{\mathrm{loc}}^{(\le \Delta_{\max})}$ be the finite truncation of the graded local operator basis from Definition~\ref{def:local_operator_spaces},
and choose the canonical flowed lifts block-by-block as organized by Theorem~\ref{thm:LaneA_associated_graded_identity} and Proposition~\ref{prop:canonical_block_uv_identity}.
Then Corollary~\ref{cor:LaneA_all_finite_truncations} implies that Lane~A applies unconditionally to every $\mathfrak A_{\mathrm{loc}}^{(\le \Delta_{\max})}$.
The resulting finite-truncation sharp-local states are compatible under the natural inclusions
\[
\mathfrak A_{\mathrm{loc}}^{(\le \Delta_{\max})}\hookrightarrow \mathfrak A_{\mathrm{loc}}^{(\le \Delta_{\max}')}
\qquad (\Delta_{\max}\le \Delta_{\max}'),
\]
and therefore define a unique state on the inductive union
\[
\mathfrak A_{\mathrm{loc}}:=\varinjlim_{\Delta_{\max}\to\infty}\mathfrak A_{\mathrm{loc}}^{(\le \Delta_{\max})}.
\]
This state extends $\omega_{\mathrm{bd}}^{(u_{\mathrm{ren}})}$ on $\mathcal A_+$ and satisfies the OS axioms on the full sharp-local algebra generated by flowed curvature composites.
\end{corollary}

\begin{proof}
Fix $\Delta_{\max}$.
By Corollary~\ref{cor:LaneA_all_finite_truncations}, the whole truncation $\mathfrak A_{\mathrm{loc}}^{(\le \Delta_{\max})}$ already satisfies both the quantitative inverse-control and the SFTE/remainder-compatibility inputs of Theorem~\ref{thm:LaneA_closure}.
Hence Theorem~\ref{thm:LaneA_closure} applies to that truncation.
If $\Delta_{\max}'\ge \Delta_{\max}$, the canonical flowed lifts and normalization conditions on the smaller truncation are the restrictions of those on the larger truncation,
so the resulting Schwinger functions agree on the overlap.
Any polynomial in sharp-local generators involves only finitely many basis elements, hence belongs to some finite truncation; define its expectation there.
Compatibility makes this independent of the chosen truncation, and the OS axioms pass to the inductive union because every OS test polynomial is finite.
The bounded-state restriction statement is inherited from Theorem~\ref{thm:LaneA_closure}(c) on each finite cap and therefore on the inductive union.
\end{proof}


\begin{theorem}[Bounded positive-time base algebra is cyclic in the full sharp-local reconstruction]\label{thm:LaneA_bounded_base_cyclic}
Fix $u_{\mathrm{ren}}\in(0,u_\ast]$ and let $\omega^{(u_{\mathrm{ren}})}$ be the full sharp-local state furnished by Corollary~\ref{cor:LaneA_full_scope}.
Let $(\mathcal H_{\mathrm{loc}},\Omega_{\mathrm{loc}},H_{\mathrm{loc}})$ be the OS reconstruction of this state from the full sharp-local positive-time algebra.
Let $\mathcal P_{+,\mathrm{sharp}}$ denote the $\ast$-algebra of finite positive-time polynomials in the renormalized local generators
$\mathcal O^{\mathrm{ren}}_j(f;\mu)$ with test functions supported in $\{x_0>0\}$.
Then:
\begin{enumerate}[label=(\roman*),leftmargin=2.3em]
\item for every $P\in\mathcal P_{+,\mathrm{sharp}}$ there exists a family $P_t^{\mathrm{bnd}}\in\mathcal A_+$ such that
\[
\|[P_t^{\mathrm{bnd}}]-[P]\|_{\mathcal H_{\mathrm{loc}}}\xrightarrow[t\downarrow 0]{}0;
\]
\item the image of $\mathcal A_+$ is dense in $\mathcal H_{\mathrm{loc}}$, equivalently $\Omega_{\mathrm{loc}}=[1]$ is cyclic for the bounded positive-time base algebra inside the full sharp-local reconstruction;
\item the centered bounded classes satisfy
\[
\overline{\mathrm{span}}\{[F-\omega^{(u_{\mathrm{ren}})}(F)]:F\in\mathcal A_+\}=\Omega_{\mathrm{loc}}^\perp.
\]
\end{enumerate}
\end{theorem}

\begin{proof}
Let $P\in\mathcal P_{+,\mathrm{sharp}}$.
Only finitely many sharp-local generators occur in $P$, so there exists an engineering cap $\Delta_{\max}$ such that
$P\in\mathfrak A_{\mathrm{loc}}^{(\le \Delta_{\max})}$.
By Corollary~\ref{cor:LaneA_all_finite_truncations}, Theorem~\ref{thm:LaneA_closure} applies unconditionally on that cap.
Choose $\varepsilon>0$ so that every test function appearing in $P$ is supported in $\{x_0\ge \varepsilon\}$.
For $t\in(0,t_\ast]$ choose $R(t)\ge 1$ with $R(t)\sqrt t\le \varepsilon/2$ and $R(t)\asymp\sqrt{\log(1/t)}$ as in Lemma~\ref{lem:RP_safe_flow_rep}.

For each sharp-local generator $\mathcal O^{\mathrm{ren}}_j(f;\mu)$ appearing in $P$, let
$\mathcal O^{\mathrm{app}}_j(t;f;\mu)$ be the centered inverse-flow approximant from \eqref{eq:inverse_flow_approx},
and let $\mathcal O^{\mathrm{bnd}}_{j,t}(f;\mu)\in\mathcal A_+$ be obtained by replacing every flowed probe inside
$\mathcal O^{\mathrm{app}}_j(t;f;\mu)$ by its RP-safe representative from Lemma~\ref{lem:RP_safe_flow_rep}.
Now define $P_t^{\mathrm{bnd}}\in\mathcal A_+$ by evaluating the same noncommutative polynomial expression as $P$
in the bounded representatives $\mathcal O^{\mathrm{bnd}}_{j,t}(f;\mu)$.
Because $P$ contains only finitely many generators, $P_t^{\mathrm{bnd}}$ is a bounded positive-time cylinder observable for each fixed $t$.

The proof of Theorem~\ref{thm:LaneA_closure}(b) already shows how to pass from convergence of the individual inverse-flow approximants
to convergence of any fixed polynomial in them: expand the finitely many monomials, apply Theorem~\ref{thm:LaneA_closure}(a)
for the SFTE/inverse-flow part, and use Lemmas~\ref{lem:RP_safe_flow_rep} and~\ref{lem:bad_event_irrelevant} to replace flowed probes by RP-safe representatives with an $o(1)$ error.
Applying that argument to the finitely many monomials appearing in $P$, in $\Theta P\cdot P$, and in the mixed products with $P_t^{\mathrm{bnd}}$, one gets
\[
\omega^{(u_{\mathrm{ren}})}(\Theta P_t^{\mathrm{bnd}}\,P_t^{\mathrm{bnd}})\to \omega^{(u_{\mathrm{ren}})}(\Theta P\,P),
\qquad
\omega^{(u_{\mathrm{ren}})}(\Theta P_t^{\mathrm{bnd}}\,P)\to \omega^{(u_{\mathrm{ren}})}(\Theta P\,P).
\]
Therefore
\begin{align*}
\|[P_t^{\mathrm{bnd}}]-[P]\|_{\mathcal H_{\mathrm{loc}}}^2
&=\omega^{(u_{\mathrm{ren}})}\!\big(\Theta(P_t^{\mathrm{bnd}}-P)(P_t^{\mathrm{bnd}}-P)\big)\\
&=\omega^{(u_{\mathrm{ren}})}(\Theta P_t^{\mathrm{bnd}}\,P_t^{\mathrm{bnd}})
-2\Re\,\omega^{(u_{\mathrm{ren}})}(\Theta P_t^{\mathrm{bnd}}\,P)
+\omega^{(u_{\mathrm{ren}})}(\Theta P\,P)
\xrightarrow[t\downarrow 0]{}0,
\end{align*}
which proves~(i).

For~(ii), by construction of the OS Hilbert space of the full sharp-local state, the span of classes $[P]$ with $P\in\mathcal P_{+,\mathrm{sharp}}$ is dense in $\mathcal H_{\mathrm{loc}}$.
Part~(i) approximates each such class by classes of elements of $\mathcal A_+$, hence the image of $\mathcal A_+$ is dense.
Since $[1]=\Omega_{\mathrm{loc}}$, this is exactly cyclicity of the bounded positive-time base algebra.
Finally, for every bounded $F\in\mathcal A_+$,
\[
[F]=[F-\omega^{(u_{\mathrm{ren}})}(F)]+\omega^{(u_{\mathrm{ren}})}(F)\,\Omega_{\mathrm{loc}}.
\]
Thus the span of centered bounded classes together with $\Omega_{\mathrm{loc}}$ is dense, and orthogonal projection to $\Omega_{\mathrm{loc}}^\perp$ gives~(iii).
\end{proof}


\begin{theorem}[Canonical scalar non-triviality witness]\label{thm:LaneA_nontriviality_witness}
Fix $u_{\mathrm{ren}}\in(0,u_\ast]$ and let $\omega^{(u_{\mathrm{ren}})}$ be the full sharp-local state from Corollary~\ref{cor:LaneA_full_scope}.
Let $\mathcal E_{\mathrm{ren}}$ denote the centered CP-even dimension-$4$ scalar field constructed from the canonical scalar block,
and let $\phi_{\ast,S}$ be the positive-time reference test function of Lemma~\ref{lem:LaneA_scalar_reference_norm}.
Then
\begin{equation}\label{eq:LaneA_nontriviality_norm_one}
\omega^{(u_{\mathrm{ren}})}\big(\Theta\mathcal E_{\mathrm{ren}}(\phi_{\ast,S})\,\mathcal E_{\mathrm{ren}}(\phi_{\ast,S})\big)=1.
\end{equation}
Consequently the OS class $[\mathcal E_{\mathrm{ren}}(\phi_{\ast,S})]$ is nonzero in $\mathcal H_{\mathrm{loc}}$,
$\mathcal E_{\mathrm{ren}}$ is not a multiple of the identity,
and the full adjoint-local sharp-local theory is non-trivial in its vacuum sector.
\end{theorem}

\begin{proof}
Fix any engineering cap containing the CP-even scalar block.
By Lemma~\ref{lem:LaneA_scalar_reference_norm}, the centered scalar component $X^S_j(\phi_{\ast,S})$ of the finite-truncation renormalized multiplet
has OS norm $1$ at every dyadic scale $t_j$.
By Theorem~\ref{thm:field_algebra_truncation}, these dyadic approximants converge in $L^2$ and hence in OS norm to the limiting centered renormalized scalar field
$\mathcal E_{\mathrm{ren}}(\phi_{\ast,S})$ inside that truncation.
Passing to the limit in the norm identity therefore yields \eqref{eq:LaneA_nontriviality_norm_one} on the finite cap.
Corollary~\ref{cor:LaneA_full_scope} and the cap-compatibility statement of Lemma~\ref{lem:LaneA_block_cap_compat}
identify the same normalized scalar field inside every larger cap and hence on the full inductive-union algebra.
Thus \eqref{eq:LaneA_nontriviality_norm_one} holds in the full sharp-local theory.
A multiple of the identity has zero centered OS class, whereas \eqref{eq:LaneA_nontriviality_norm_one} shows that
$[\mathcal E_{\mathrm{ren}}(\phi_{\ast,S})]$ has norm one.
Hence the field is not a multiple of the identity and the vacuum sector is non-trivial.
\end{proof}

\begin{remark}[Why the inductive-limit step is only a compatibility argument]\label{rem:LaneA_inductive_compatibility}
The only genuinely new bookkeeping input in Corollary~\ref{cor:LaneA_full_scope} is the coherent choice of basis and normalization.
At each exact-dimension symmetry block, the local basis and flowed basis are written using the same complete-contraction formulas,
and the normalization moments are fixed once and for all.
When a larger engineering cap is introduced, Lemma~\ref{lem:coherent_Z} ensures that the fields already constructed in smaller caps are literally the same fields inside the larger cap.
Hence the inductive-limit passage is a compatibility argument, not a second renormalization problem.
\end{remark}

\begin{remark}[Lane~A interface status]
The role split for Packets~5--7 is frozen in Table~\ref{tab:LaneA_packet_interface_map} of Appendix~\ref{app:normalization_crosscheck}.
Theorem~\ref{thm:finite_truncation_inverse_control} and Corollary~\ref{cor:finite_truncation_SFTE} provide only the unconditional finite-truncation inverse / remainder data;
Theorem~\ref{thm:finite_cap_flowed_sector_bridge} is the first mixed-correlator closure node;
Theorem~\ref{thm:LaneA_closure} and Corollary~\ref{cor:LaneA_full_scope} perform the capwise extension and then the inductive-union upgrade;
and Theorem~\ref{thm:LaneA_bounded_base_cyclic} is the first theorem that adds bounded-base cyclicity.
No later sentence should be read as re-proving that packet chain in prose.
\end{remark}
